% ================================================================== 1
\section{Resumen}
\label{sec:resumen}

RSA se construye sobre tres problemas de teoría de números: decidir si un
número es primo (para generar las claves), calcular potencias modulares
grandes (para cifrar y descifrar) y factorizar enteros (el ataque del que
depende su seguridad). Cada uno admite algoritmos con costos muy distintos. El
proyecto desarrolla una biblioteca educativa en Python, \texttt{rsa\_edu}, que
implementa desde cero, para cada etapa del pipeline RSA, una alternativa base y
una alternativa optimizada, y mide si la optimizada hace el sistema más rápido
sin cambiar el resultado.

Se comparan la exponenciación binaria con la ventana deslizante; la división
de prueba con Miller--Rabin (aleatorio y determinista); la generación de primos
con y sin criba de primos pequeños; el descifrado directo con el Teorema Chino
del Resto; y, para el ataque, la división de prueba con Pollard Rho con
detección de ciclos de Floyd y de Brent. Las claves siguen FIPS~186-5 y las
primitivas de cifrado PKCS~\#1~v2.2. La corrección se verifica con 397 pruebas
automatizadas (94\,\% de cobertura) contra referencias externas (\texttt{pow}
de Python y \texttt{sympy}), y la eficiencia con un diseño experimental
reproducible (semilla fija) que mide tiempo, operaciones y memoria.

Las mediciones coinciden con el análisis de complejidad: la división de prueba
crece como $O(\sqrt n)$ y Pollard Rho como $O(n^{1/4})$; la ventana reduce las
multiplicaciones de $1{,}5$ a $1{,}18$ por bit del exponente; la criba evita el
83\,\% de las pruebas de Miller--Rabin (la teoría predice 85\,\%); y el
descifrado por TCR es de 2,3 a 3,3 veces más rápido. Con la misma clave y los
mismos cifrados, el pipeline RSA optimizado es 3,1 veces más rápido con claves
de 512 bits y 5,8 veces con 2048 bits. Optimizar la generación de primos y el
descifrado reduce el costo total de RSA entre 3 y 6 veces, y la mejora crece
con el tamaño de la clave; el cifrado con $e = 65537$ ya es barato y no se
beneficia.
